\documentclass[11pt]{article}

% === Page Layout ===
\usepackage[a4paper, margin=1in]{geometry}

% === Packages ===
\usepackage{graphicx}
\usepackage{comment}
\usepackage{fancyhdr}
\usepackage{enumitem}
\usepackage{amsmath}

% === Header/Footer Setup ===
\pagestyle{fancy}
\fancyhf{}
\fancyhead[L]{Multimodal Deep Learning}
\fancyhead[R]{Answer Sheet 2}
\fancyfoot[C]{\thepage}

\begin{document}

% === Title Information ===
\begin{center}
    {\Large \textbf{Multimodal Deep Learning}} \\ [0.3em]
    {\large \textbf{Answer Sheet 2}} \\ [0.3em]
    \large Date: May 6, 2025
\end{center}

\vspace{1cm}

\subsection*{Exercise 1: Numerical Features: Normalization vs. Standardization}

\noindent
A one-dimensional sensor produces readings

\[
    \texttt{x} = [2.7, \, 3.1, \, 9.4, \, 6.5, \, 4.2, \, 3.8].
\]

\begin{enumerate}[start=3, label=(\alph*)]
    \item {
        When would you prefer Min–Max scaling over z-scaling, and vice-versa? \\
        \textbf{Answer: } \textbf{Min-Max scaling} is preferred when the true scale of information are sensitive to the quality of the output.
        While \textbf{z-scaling} is preferred when the input is approximately Gaussian, and the distribution of features can enhance the quality of the output
        }
    \item {
        \textbf{Binning / Discretization:} Why convert a continuous feature (e.g. ``age") into bins (``young", ``middle-aged", ``senior")? Give one benefit and one drawback. \\
        \textbf{Answer: } One good benefit is discretization can help group a continuous feature into a simpler format which also help reduce the memory consumption, 
        while one drawback is its lack of feature awareness and ruin system quality, particularly when binning algorithm is too naive or trying to do without sufficient data understanding.
        }
\end{enumerate}

\subsection*{Exercise 2: Categorical Features}

\noindent
You have a categorical attribute \texttt{Color} with values

\begin{center}
    [Red, Blue, Green, Blue].
\end{center}

\begin{enumerate}[label=(\alph*)]
    \item {
        Show the One-Hot encoded matrix. \\ 
        \textbf{Answer: }
        \[
            \texttt{X} = 
            \begin{bmatrix}
            0 & 0 & 1 \\
            1 & 0 & 0 \\
            0 & 1 & 0 \\
            1 & 0 & 0 \\
            \end{bmatrix}
        \]
        }
    \item {
        Show the Label-Encoding (integers assigned alphabetically from 0). \\ 
        \textbf{Answer: }
        \[
            \text{Blue} \mapsto 0, \ \text{Green} \mapsto 1, \ \text{Red} \mapsto 2
        \]
        }
    \item {
        What can go wrong if you feed the labels from (b) directly into a linear or distance-based model? \\ 
        \textbf{Answer: } The model will not be able to achieve its best performance because the label-encoding method does not manifest the exact relationship between each color.
        However, One-Hot encoding is a better option because it presents the vectors of color in vector space avoiding assumption of order or magnitude.
        }

\end{enumerate}

\subsection*{Exercise 3: Text Features}

\noindent
Byte-pair encoding (BPE) corpus

\begin{center}
    [``low'', ``lower'', ``lowest'', ``newest'']
\end{center}

\noindent
Perform 4 BPE merges. Show the merge pairs and the updated vocabulary after each step.

\noindent
\textbf{Answer: } Please see the answer provided by the lecturer.

\subsection*{Exercise 4: Johnson–Lindenstrauss Lemma Bounds}

\noindent
For simplicity in this exercise, we use the bound C = 8.

\begin{enumerate}[label=(\alph*)]
    \item {
        \(N\) = \(10^6\), \(\epsilon\) = 0.1: compute the minimum \(k\).
        \[
            k \geq \frac{C \log N}{\epsilon^2} = \frac{8 \log 10^6}{0.1^2} \approx 11052.
        \]
        Minimum \(k\) = 4800.
        }
    \item {
        With \(k\) = 500 and \(N\) = \(10^6\), what is the best guaranteed distortion \(\epsilon\)?
        \[
            \epsilon \geq \sqrt{\frac{C \log N}{\epsilon^2}} = \sqrt{\frac{8 \log 10^6}{500}} \approx 0.47.
        \]
        The best guaranteed distortion \(\epsilon\) = 0.47 = 47\%.
        }
\end{enumerate}

\subsection*{Exercise 5: Sparse vs. Dense Distance Functions}

\begin{enumerate}[label=(\alph*)]
    \item {
        \textbf{Sparse binary vectors: } \\ 
        \textbf{Answer: } \textbf{Euclidean distance} cannot be utilized effectively with sparse data because the \texttt{0} and \texttt{1} values do not carry meaningful information for calculating distance, unlike in dense data, where features tend to have more varied and informative values.
        }
    \item {
        \textbf{Unit-normalized dense vectors:} \\
        \textbf{Answer: } From the provided answer, ranking values of \(s \mapsto d(u,v) = \sqrt{2-2s}\) by both of them are equivalent after \(L2\)-norm.
        This can refer that for certain tasks which require the order of values, we can effectively approximate the sequence of data using Euclidean distance -- lower computational cost (e.g., Approximate Nearest Neighbor (ANN) Search)
        }
    \item {
        \textbf{Dense count vectors of different length: }
        \textbf{Answer: } Cosine similarity is invariant to the length of the document, unlike Euclidean distance.
        }
\end{enumerate}
\end{document}